\subsection{Example 3}\label{sec:interval-example2-1}

\subsubsection{Problem description}\label{subsec:interval-example2-1-problem}

A regional economic survey asks 200 small-business owners whether they expect profits to grow next quarter. The goal is to estimate the population proportion who are optimistic about growth.

Sample results: 124 owners answered ``yes.''

Construct a 95\% confidence interval for the population proportion of optimistic owners.


\subsubsection{(C7) Compute the sample proportion, standard error, and confidence interval}\label{subsubsec:interval-example2-1-c7}

The binary variable is whether an owner expects profits to grow (yes = success, no = failure). The sample proportion is

\[
\hat{p} = \frac{124}{200} = 0.62
\]

Because proportions come from binary outcomes, variability depends on \(\hat{p}\) itself:

\[
SE_{\hat{p}} = \sqrt{\frac{\hat{p}(1 - \hat{p})}{n}} = \sqrt{\frac{0.62(0.38)}{200}} \approx 0.034
\]

This shows that larger samples shrink the standard error and yield tighter intervals. For 95\% confidence, use \(z_{\alpha/2} = 1.96\):

\[
E = 1.96 \cdot 0.034 \approx 0.067
\]
\[
\hat{p} \pm E = 0.62 \pm 0.067 = (0.553, 0.687)
\]

The interval gives the plausible range of the true population proportion.

\subsubsection{(C4) Interpret the interval in economic or financial terms}\label{subsubsec:interval-example2-1-c4}

With 95\% confidence, the true share of optimistic businesses lies between about 55.3\% and 68.7\%. Policy planners can use this range to estimate how many firms may seek expansion loans or increase hiring, aligning support programs with realistic demand.

\subsubsection{(C10) Conclude about a statistical parameter in finance and economy}\label{subsubsec:interval-example2-1-c10}

The survey supports concluding that the population proportion of optimistic small-business owners is roughly 0.55--0.69, with a central estimate around 0.62. Economic planners can use this parameter to forecast loan demand and design support programs that match expected business sentiment.

\vspace{6pt}
\noindent\hyperref[table-of-contents]{Back to Top}
